% ============================================================
% VII. KẾT LUẬN
% ============================================================
\section{Kết luận}

\begin{frame}{Kết luận: các điểm đã chốt}
  \begin{block}{Đóng góp đã hoàn thành}
    \begin{itemize}
      \item \textbf{Phân tích}: long-tail imbalance, texture confusion, domain gap, boundary yếu trên FoodSeg103
      \item \textbf{Baseline}: BiSeNet được chuẩn hóa và đánh giá trên \textbf{FoodSeg103}
      \item \textbf{Mở rộng}: tái cấu trúc nhãn 103 $\to$ 75 được tách thành nhánh thực nghiệm bổ sung
      \item \textbf{Phân tích lỗi}: xác định sink class, nhầm lẫn có cấu trúc, định vị vùng yếu
    \end{itemize}
  \end{block}
\end{frame}

\begin{frame}{Kết luận: hướng tiếp theo}
  \begin{alertblock}{Hướng tiếp theo}
    Kiểm tra hệ thống 3 module đề xuất: \textbf{Texture Branch + Relation GNN + Boundary Supervision} trên \texttt{FoodSeg103} nguyên thủy, đồng thời mở rộng theo hướng \textbf{Curriculum Learning}
  \end{alertblock}
  \vspace{0.2em}
  \begin{block}{Kỳ vọng khi đánh giá vòng sau}
    \begin{itemize}
      \item tăng mIoU ở lớp hiếm và lớp có texture tương tự
      \item giảm hiện tượng sink class trong confusion matrix
      \item cải thiện chất lượng biên mà vẫn giữ chi phí nhẹ
      \item kiểm chứng hiệu xuất của curriculum/self-paced learning cho tập dữ liệu mất cân bằng Foodseg 103.
    \end{itemize}
  \end{block}

  \textbf{Mã nguồn:} \url{https://github.com/ariesanhthu/SemanticSegmentation-FoodSeg103}
\end{frame}
